% ---
% titre: Prise de contrôle de la table de multiplications
% theme: NO
% chapitre: Nombres naturels et décimaux
% sous_chapitre: Calcul réfléchi et mental
% niveau: [9e]
% type: JEU
% tags: [c-livrets,f-groupe,f-station,f-concours,o-des,o-sablier,p-entrainement,p-differentiation]
% difficulte: 1
% corrige: false
% source: pdf
%source_originale: 00-SourcesLaTeX/01-NO/JEU_NO-NbNaturelsDecimauxProcedureCalculReflechi_PriseControleTableMutliplication
%notes: Trouvé sur TPT, créé par the elementary math consultant (times table takeover); In this game, students do lots of multiplication! The objective of the game is to “takeover” the multiplication table. Player 1 rolls two dice and uses the sum as her first factor. She does the same to get her second factor. She uses the multiplication chart to find the product and uses her color to color it in. Then she finds all the other multiplication facts with the same product and colors those in too. Finally, she completes the multiplication fact sheet for each multiplication problem she colored in. Player 2 follows the same procedure. If he happens to get the same product, he tries to find products on the multiplication table that Player 1 missed, and if there are none he loses a turn! Play ends when class is over, the board is all colored in (besides the ones facts because you can’t roll a 1 with two dice), or when over 50% of the board is colored in by one player (the teacher can choose which option best fits for her class). Play a round or two with the class to show them how it works, and then split them into pairs or small groups to play. About ten minutes before class ends, gather the group and use the discussion questions to guide a reflective conversation.
%notes: Ask students to list all the multiplication problems that have a product of 12 (1x12, 2x 6, 3x4, etc), 24, and 36. Does anyone see any patterns or connections? Can someone share some strategies for remembering multiplication facts? How can I solve problems that are times 12? (times ten plus times 2) A prime number is a number that only has a factor of one and itself. Did anyone have a prime number for a product? What products were the best to get for taking over more of the table? Why do those composite numbers have so many factors? How can what you did today help you solve division problems?
% ---